\chapter{Programming the Device}
\label{ch:programming}

This chapter collects what a program does in what order. It adds no new rules; each rule is stated where the feature is
described, and the sections of the earlier chapters are cited.

\section{Starting Up}
\label{sec:startup}

At reset the picture is blank (\texttt{GD} is 1), every writable register is zero except \texttt{CONFIG}
(Section~\ref{sec:reset}), and video memory, the palette and the sprite table are undefined. A program starts the
device in this order:

\begin{enumerate}
	\item Install the handler for the vertical sync interrupt, \activelow{NMI}, before anything enables it
	      (Section~\ref{sec:interrupts}).
	\item With the picture still blank, write the palette (\texttt{0x6000}--\texttt{0x60FF}), the tile data and the tilemaps into
	      video memory (Chapter~\ref{ch:memory-map}). Video memory is open to the CPU while \texttt{GD} is 1.
	\item Fill the sprite table, hiding every entry the program does not use by setting its layer field to 3
	      (Section~\ref{sec:sprite-entry}).
	\item Write the tilemap address, tile data address, palette configuration and scroll registers of the backgrounds in use, the
	      \texttt{STILES} register, and the \texttt{TILESIZE} register.
	\item Write \texttt{INTCTL} to enable the interrupts the program uses (\texttt{VIE}, and \texttt{LIE} with \texttt{LINE} for
	      a line-compare interrupt).
	\item Write \texttt{CONFIG} with the mode, the brightness and \texttt{GD} = 0. The picture appears from the next line.
\end{enumerate}

\section{One Frame}
\label{sec:frame-flow}

Figure~\ref{fig:frame-flow} shows a frame as the program sees it. The device draws lines 0 to 223 without help; line 0
is always blank (Section~\ref{sec:first-line}). At the first cycle of line 224 it raises the \activelow{NMI} (if
\texttt{VIE} is set) and sets the \texttt{VBL} bit of the Status register, and the handler has the rest of vertical
blanking, 38 lines or 25,916 cycles on NTSC and 88 lines on PAL, to change video memory, the sprite table and the
registers for the next frame.

\begin{figure}[H]
	\centering
	\begin{tikzpicture}[
			box/.style={draw, rounded corners=2pt, minimum height=0.8cm, text width=3.6cm, align=center, font=\scriptsize},
			dev/.style={box, fill=sirclightgray},
			flow/.style={-{Stealth[length=2mm]}, thick}
		]
		\node[dev] (draw) at (0,0) {Lines 0 to 223: the device draws the picture};
		\node[dev] (nmi) at (0,-1.6) {Line 224: \activelow{NMI} (if \texttt{VIE})};
		\node[box] (h1) at (5,-1.6) {Handler: move sprites, write scroll registers, update tilemaps and tiles};
		\node[box] (h2) at (5,-3.2) {Return from the interrupt};
		\node[dev] (blank) at (0,-3.2) {Lines 225 to the end of the frame: no drawing};
		\draw[flow] (draw) -- (nmi);
		\draw[flow] (nmi) -- (h1);
		\draw[flow] (h1) -- (h2);
		\draw[flow] (h2) -- (blank);
		\draw[flow] (blank.south) -- ++(0,-0.6) -| (-2.6,0) -- (draw.west);
	\end{tikzpicture}
	\caption{A Frame}
	\label{fig:frame-flow}
\end{figure}

A program that does not use the interrupt waits for the \texttt{VBL} bit instead: it reads \texttt{STATUS} until bit 7
is set, and it knows how far into the frame it is from \texttt{VLINE}. A program that polls must finish before
\texttt{VBL} falls at line 0, so it reads \texttt{VBL} again before each write that matters near the end.

A \texttt{DMAT} of 255 words takes 522 cycles (Section~\ref{sec:cpu-clock}), so a program doing nothing else could move
at most about 12,600 words in an NTSC vertical blanking interval. A program that needs more than it can move in one
frame spreads the work over several frames, or turns the picture off (\texttt{GD} = 1) for the time it needs.

\section{Programming Cautions}
\label{sec:cautions}

\begin{enumerate}
	\item \textbf{Video memory.} A write to video memory or to the sprite table outside vertical blanking is lost, and a
	      read returns zero, unless \texttt{GD} is 1 (Section~\ref{sec:memory-vram-access}). A program knows video memory is open when \texttt{VBL} in the
	      Status register is set. The palette and the registers can be written at any time.
	\item \textbf{The sprite table.} The table is undefined at reset. A program must hide the entries it does not use before
	      it clears \texttt{GD}, or random sprites appear.
	\item \textbf{NMI.} The \activelow{NMI} cannot be masked. Install the handler before setting \texttt{VIE}. A
	      \texttt{DMAT} that is running at line 224 delays the interrupt until it ends (Section~\ref{sec:interrupts}).
	\item \textbf{Line compare.} The \activelow{IRQ4} line stays asserted until the program writes \texttt{LINE}. A handler
	      that does not write \texttt{LINE}, or clear \texttt{LIE}, is entered again as soon as it returns.
	\item \textbf{Registers during the picture.} The device samples the scroll registers once per line, and a change made during a line takes effect on the next line. It reads the tilemap address, tile data address and palette configuration registers as it fetches, so a change during the visible part of a line can mix old and new settings within that line: change them only in vertical blanking or while \texttt{GD} is 1 (Chapter~\ref{ch:mode2}).
	\item \textbf{Changing mode.} Change \texttt{MODE} only in vertical blanking or while \texttt{GD} is 1. The values 3 to 7
	      of \texttt{MODE} are reserved and must not be written.
	\item \textbf{Reserved registers.} Software must not write the reserved registers and reserved bits
	      (Chapter~\ref{ch:registers}). A read of an open bus offset returns zero.
	\item \textbf{Scroll registers.} The scroll registers are 10 bits. A program keeps its position in a wider variable and
	      writes the low 10 bits; in Mode 2 it also adjusts \texttt{SCRY} when it brings \texttt{SCRX} back past 1024
	      (Section~\ref{sec:m2-moving}).
	\item \textbf{Reading video memory.} A read of video memory takes one wait state, and two straight after a write, and it returns zero outside vertical blanking and forced blank (Section~\ref{sec:signal-vram}). A program keeps in RAM what it will need again, instead of reading it back.
	\item \textbf{Sprites per line.} A line shows at most 32 sprites; the rest are dropped and \texttt{SOV} is set. Read
	      \texttt{SOV} in vertical blanking (Section~\ref{sec:sprite-budget}).
	\item \textbf{Transparent pixels.} Pixel value 0 of a background or a sprite is transparent. If every layer is transparent
	      at a pixel, the device shows palette entry 0 (Section~\ref{sec:compositing}).
	\item \textbf{PAL and NTSC.} A program that depends on the number of lines or the frame rate reads \texttt{OUT} in the
	      Status register.
\end{enumerate}
